\cotefile{29}
\scanpage{112}
Notes pour SGA 7
% archive sheet number: 111
\scanpage{117}
Comportement fonctoriel des groupes d'inertie.

\handwritten{1.} Soient $X$ un schéma, $U$ un ouvert de $X$, $\handwritten{G}$ un groupe profini, $U'$ un
revêtement principal de $U$ de groupe $G\deleted{\unsure{'}}$, $x$ un point \inserted{\handwritten{\unsure{géom|géom.}}} de $X$, $\bar X$ le localisé
strict de $X$ en $x$, $\bar U$ l'image inverse de $U$ dans $\bar X$, $\bar U'=U'\times_U\bar U$. On suppose
donné de plus un schéma $\bar X'$ sur $\bar X$, avec des opérations de $G$ sur $\bar X'$, de
telle façon que le \deleted{\unsure{espace|space|l'espace}}$\inserted{\bar U\text{-schéma}}$ à opérateurs $\bar X'\handwritten{\mid}\bar U$ co\deleted{\unsure{p}}ïncide avec $\bar U'$.

Pratiquement, on sera intéressé par le cas où $x\handwritten{\in}X-U$ et où $x$ est
adhérent à $U$, où $\bar X'$ est entier sur $\bar X$ et où $\underline O_{\bar X}=\bar p(\underline O_{\overline{X'}})^G$, $\bar p:\overline{X'}\handwritten{\to}\bar X$ étant
la projection canonique, - de sorte que $G$ opère transitiv\deleted{\unsure{i|x}}ement sur
l'ensembl\unsure{e|é} fibre $\overline{X'}_{\bar x}$. De plus, on \deleted{\unsure{se place|se plac}} sera \inserted{\handwritten{le plus souvent}} \deleted{\unsure{d}}sous des conditions où
on peut construire canoniquement $\bar X'$ en termes de $\bar U'$; par exemple, si on
suppose $\deleted{\unsure{O|U}}X$ normal en $x$, donc $\bar X\deleted{\unsure{'}}$ normal, et $\bar U$ non vide, on définira $\bar X'$ comme
le normalisé de $\bar U$ relativement à $\bar U'$ sur $\bar U$ (i.e. comme limite projective
des nor\deleted{\unsure{al|ma}}malisés d\deleted{\unsure{es}}ans les $\overline{U'(i)}$ sur $\bar U$, $U'$ \deleted{\unsure{u}}étant é\deleted{\unsure{e}}crit comme limite projective de schémas principaux $U'(i)$ sur les groupes quotients finis $G(i)$ de
$G$).

Si maintenant \deleted{\unsure{$x$ est}} $\handwritten{\xi}$ est un point de $\overline{X'}_{\bar x}$, on peut considérer le
groupe de stabilité $G_{\handwritten{\xi}}$ de $\handwritten{\xi}$, formé des $g\in G$ tels que $g.\handwritten{\xi}=\handwritten{\xi}$; ce groupe
prendra le nom de \uline{groupe d'inertie} \uline{relatif à} $\handwritten{\xi}$. Dans le cas envisagé plus
haut où les points de $\bar X'$ sur $\bar x$ sont tous conjugués sous $G$, on voit que
pour $\handwritten{\xi}$ variable\deleted{\unsure{s}}, les $G_{\handwritten{\xi}}$ parcourent une classe de sous-groupes conjugués
de $G$, qu'on appelera \inserted{(par abus de langage)} \uline{les sous-groupes d'inertie} \deleted{de} \uline{relatifs à} $x$. Noter
qu'en toute rigueur cette classe de sous-groupes \deleted{d'inertie} n'est pas
déterminée par la seule connaissance de $x\handwritten{\in}X$, mais dépend du choix du
prolongement de $\overline{U'}$ en $\overline{X'}$. Lorsque $X$ est normal en $x$ et que $x$ est spécialisation d'un point de $U$ (ce qui signifie que $x$ est adhérent à $U$ si $\mathrm{esp}\handwritten{(}X\handwritten{)}$
est localement noethérien), il sera sous-entendu, \deleted{\unsure{-}}sauf mention du contraire, que le choix en question est fait par normalisation, comme expliqué
plus haut.
% archive sheet number: 116
\scanpage{118}
% The writing on this face is upside down relative to the archive foliation.
\mspage{\handwritten{1}}
\deleted{\[
\begin{aligned}
&{}_{\underline a}P\ \unsure{\Omega|\otimes|\cap}\ {}_{\underline b}Q
\end{aligned}\tag*{(2.1.14.6)}
\]}
% archive sheet number: 117
\scanpage{124}
``\uline{Vrai}'' groupe fondamental

[feuilles \unsure{emboîtées|emboitées|inutiles|anciennes}]
% archive sheet number: 123
\scanpage{126}
\mspage{\handwritten{2}}

de schémas en groupes finis et plats sur $S$
\[
e\handwritten{\to}H_{\mathrm{inf}}\handwritten{\to}H\handwritten{\to}H_{\text{sép}}\handwritten{\to}e
\]
où $H_{\mathrm{inf}}$ est purement infinitésimal (i.e. la projection $H_{\mathrm{inf}}\handwritten{\to}S$
est géométriquement injectif) et où $H_{\text{sép}}$ est non ramifié sur $S$.
(Dans le cas d'un corps de base $k$, une telle suite exacte est
déduite d'une suite exacte analogue pour un groupe algébrique
fini \inserted{$G$} sur $k$; d'ailleurs si $k$ est parfait, cette suite splitte canoniquement, car $G_{\text{réd}}$ est \inserted{alors} un sous-groupe de $G$, isomorphe à $G_{\text{sép}}$ par la
projection $G\quad G_{\text{sép}}$. Faire attention cependant que $G_{\text{réd}}$ va \unsure{opérer|opérser}
sur $G_{\mathrm{inf}}$, on aura seulement un produit semi-direct). Enfin

\deleted{\unsure{La condition (v) est certainement provisoire}}

(vi) $S$ est réduit.

Les conditions (v) et (vi) ont une sale gueule, et sont essentiellement provisoires. Elles servent à \deleted{assurer la validité} \inserted{\handwritten{cause}} du

\uline{Lemme} Soit $H$ comme dans (v), $S$ comme dans (vi), $P$ un fibré principal homogène sous $H$, $Q$ un autre, $u,v$ deux isomorphismes de $P$ dans $Q$
transformant un \handwritten{``}point \deleted{\unsure{rationnel}} \inserted{marqué}\handwritten{''} donné en un autre donné, alors
$u=v$. (On est ramené, en tordant $H$, au cas où $P$ est trivial, et
alors cela résulte du fait suivant: une section de $Q$ est un isomorphisme de $S$ sur une composante connexe de $Q_{\text{réd}}$).

Remarque néanmoins que les conditions (i) à (vi) n'excluent
pas des groupes structuraux tordus (relativement à un corps de base
$k$).

Soit maintenant $a$ un ``point marqué'' de $S$ (i.e. une extension
algébriquement close du corps résiduel d'un $s\handwritten{\in}S$). Pour tout $G\handwritten{\in}\underline G$,
on désigne par $Z(S,a;G)$ ou simplement $Z(G)$ l'ensemble des classes
(à isomorphisme près) de fibrés principaux homogènes sur $S$, de
groupe $F(G)$, munis d'un point marqué au dessus de $a$. Evidemment
$Z(G)$ est un foncteur de $\underline G$ dans la catégorie des ensembles. Ce
foncteur satisfait aux conditions suivantes: 1) Il commute aux
produits (\deleted{\unsure{trivial), et|trivial}} car $F$ y commute) 2) il commute aux noyaux
de couples, en d'autres termes si $G''\handwritten{\to}G\handwritten{\rightrightarrows}G'$ est exact dans $\underline G$,
alors $Z(G'')\handwritten{\to}Z(G)\handwritten{\rightrightarrows}Z(G')$ est exacte, i.e. $Z(G'')$ s'identifie
à l'ensemble des éléments de $Z(G)$ dont l'image dans $Z(G')$ par $Z(u)$
\deleted{\unsure{(}}et $Z(v\deleted{\unsure{a}})$ est la même. (Ceci résulte de l'exactitu\unsure{de|é} de $F$, et du
lemme). \unsure{t|c}
% archive sheet number: 125
\scanpage{129}
\mspage{\handwritten{5}}

Lorsqu'il n'y a pas de corps de base, il n'est pas clair cependant
comment et à quels groupes fondamentaux \deleted{\unsure{on peut|on puisse}} pour $X,Y$ on peut
rattacher \deleted{par une suite} les groupes \unsure{fondementaux|fondamentaux} des fibres,
pour tenir lieu de suite exacte d'homotopie. S'il y a un corps de
base, une \deleted{\unsure{hypothèse}} \inserted{conjecture} de départ serait que \deleted{l'image inverse par $f$
de} $\handwritten{\underline\pi}_1(X/Y)$, le pro-schéma en groupes \inserted{$\handwritten{\underline{\pi}_1(X)}$} sur $X$ qui définit le système
local des groupes fondamentaux \deleted{\unsure{aux po}} de $X$ en ses divers points,
et enfin l'image inverse par $f$ du pro-schéma en groupes analogue $\handwritten{\underline{\pi}_1(Y)}$
sur $Y$, sont reliés par une suite exacte \deleted{(qui serait définie d'ailleurs
indépendamment de l'existence d'une section de $X$ sur $Y$, car le
premier \deleted{\unsure{groupe}} pro-schéma en groupes \inserted{sur $X$} nommé serait défini en tous
cas).} \deleted{\unsure{Si $Y$ est le spectre d'un anneau local complet, notamment un
anneau de valuation discrète, on devrait pouvoir mettre|Si $Y$ est le spectre d'un anneau local complet, notamment un anneau de valuation discrète, on devrait pouvoir s'y mettre.}} Il y a de
la joie en perspective pour les groupes d'homotopie supérieurs \ldots

Il n'est pas dit que la conjecture envisagée soit plus difficile
à démontrer que la partie déjà connue pour les groupes fondamentaux
ordinaires. Elle impliquerait que pour des schémas complets $X,Y$ \deleted{\unsure{q}}au
dessus d'un corps alg clos $k$, on a $\handwritten{\pi}_1(X\times Y/k)=\handwritten{\pi}_1(X/k)\times\handwritten{\pi}_1(Y/k)$.
Utilisant tes raisonnements, on en conclut que si $X$ est une
variété abélienne sur $k$, alors $\handwritten{\pi}_1(X/k)$ est abélien, et qu'un
revêtement principal minimal $X'$ de $X$ est une variété abélienne isogène à $X$. On en déduit
\[
\handwritten{\pi}_1(X/k)=\underset{\handwritten{\longleftarrow}}{\lim}\ {}_nX
\]
où ${}_nX$ est le noyau (avec sa partie infinitésimale aussi, bien sûr)
de la multiplication par $n$, l'homomorphisme de ${}_{mn}X$ dans ${}_nX$ étant
induit par la multiplication par $m$. (Utilisant Cartier, ce groupe
fondamental est dual, au sens de Cartier, au groupe ind-algébrique
limite inductive des ${}_nX^{\handwritten{*}}$, où $X^{\handwritten{*}}$ est la variété duale de $X$).
Prenant la $p$-composante de ce groupe fondamental, on trouve ce qui,
pour le nombre premier $p$, doit jouer le rôle du module de Weil. Il
est hors de doute (et Cartier doit l'avoir fait) qu'il est possible
d'associer \deleted{à un} fonctoriellement à un groupe proalgébrique infinitésimal abélien un module sur l'anneau de Witt, et \deleted{\unsure{de vé}} qu'il \deleted{\unsure{soi}}\handwritten{soit}
facile de vérifier qu'en l'occurence ce module se décompose en les
trois morceaux que tu connais bien (correspondants aux trois types
principaux de \inserted{$\handwritten{p}$-}groupes algébriques abéliens), \deleted{qui} On trouve ainsi
% archive sheet number: 128
\scanpage{130}
\mspage{\handwritten{6}}

pour ton théorème une formulation plus naturelle (qui devrait \deleted{\unsure{su}}
fournir une démonstration uniforme, sans distinction du cas $\handwritten{\ell\ne}p$
et $\handwritten{\ell}=p$), de même qu'on voit en même temps que ta somme abracadabrante se comporte bien quand on fait varier $X$ dans une famille, i.e
si $X$ est \deleted{\unsure{en fait un}} un schéma abélien sur $Y$ (car alors les ${}_nX$
seront d'excellents schémas en groupes finis et plats sur $Y$, dont on
peut prendre formellement la limite projective \ldots).

J'espère arriver dans l'année prochaine à une théorie satisfaisante du groupe fondamental, \deleted{\unsure{ce qui}} et achever la rédaction des
Chapitres IV,V,VI,VII (ce dernier étant le groupe fondamental), en
même temps que des catégories. Dans deux ans résidus, dualité,
intersections, Chern, Riemann-Roch. Dans trois ans cohomologie de
Weil, et un peu d'homotopie si Dieu veut. Et entre-temps, je ne sai\deleted{\unsure{t}}\unsure{s|t}
quand, le grand théorème d'existence avec Picard etc, un peu de
courbes algébriques, \inserted{les schémas \unsure{abéliens|abélians}.} Sauf difficultés imprévues ou enlisement,
le multiplodoque devrait être fini d'ici 3 ans, ou 4 au maximum. On
pourra commencer à faire de la géométrie algébrique !

Bien à toi\\
% archive sheet number: 129
\scanpage{131}
\diagram{\begin{tikzcd}[ampersand replacement=\&,row sep=small]
X \arrow[d,"f"]\\ S
\end{tikzcd}}
muni d'une section $\sigma$, $f$ propre plat, $f_*(\mathcal O_X)=\mathcal O_S$ \uline{\unsure{universellt|universellement}}.

$G$ schéma en groupes \deleted{\illegible{1}} fini et \deleted{séparable} \inserted{étale} sur $S$
[le cas important est celui où $G=S\times\Gamma$].

$F_G(T)=$ ens des classes à isom près des
fibrés principaux homogènes sur $X_T$, de groupe
$G_T\times_T X_T=(G\times_{S} X_T\deleted{\unsure{\times X|X}}$, \uline{munis} d'une section
marquée le long de $\sigma_T$, i.e. \inserted{d'un isom} $\sigma_T^*(X')\simeq G_T$.

N.B. Tout automorphisme d'un tel torseur
est trivial, \unsure{grâce au|grâce au seul} fait que $G$
est \deleted{\unsure{sép}} étale sur $S$. Il en serait de
$=$ pour $G$ commutatif. Cas général??

Le foncteur $F_G(T)$ en $T$ est-il
représentable ?? Si oui, le \unsure{Sch|$H$} sur $S$
qui le représente est \uline{étale} sur $S$, et
\uline{quasi fini sur} $S$ [car à fibres
\uline{finies}].

% margin note written upside down; the two long strokes cross the note and diagram
\margin{\deleted{4.9, \unsure{4.12|4.10|4.11}.) \quad 4.9.}}
\diagram{\begin{tikzcd}[ampersand replacement=\&,column sep=small,row sep=small]
\& \deleted{C}\arrow[r]\arrow[d,dash] \& \deleted{\overset{\cdot}{C'}}\arrow[r,dash,shift left=1]\arrow[r,dash,shift right=1]\arrow[d] \& \deleted{C''}\arrow[r,dash,shift left=2]\arrow[r,dash]\arrow[r,dash,shift right=2]\arrow[d,dash] \& \deleted{C'''}\arrow[d,dash]\\
\deleted{\text{\uuline{\unsure{épuise|épim.}}}} \& \deleted{\underset{\cdot}{C_0}}\arrow[r,dash] \& \deleted{\underset{\cdot}{C'_0}}\arrow[r,dash,shift left=1]\arrow[r,dash,shift right=1] \& \deleted{C''_0}\arrow[r,dash,shift left=2]\arrow[r,dash]\arrow[r,dash,shift right=2] \& \deleted{C'''_0}\\
\&\& \deleted{\underline{\text{pl. fid}}} \& \deleted{\underline{\text{plat}}} \&\\
\&\& \& \deleted{\underbrace{\hspace{6em}}_{\text{\unsure{exact|exacts}}}} \&
\end{tikzcd}}

Noter que le foncteur $F_G$
est \deleted{\unsure{compati|compatible}} exact à \deleted{\unsure{g}}gauche
pour les morphismes fid. plats
et quasi-compacts.

\deleted{\illegible{2}}
\[
\deleted{G\quad H^1(X,G\quad G\quad G\ \underline{\unsure{\text{tordu}}}}
\]
% archive sheet number: 130
\scanpage{132}
% Printed orientation is upside down relative to the archive foliation.
\mspage{I-6}

subset of $\underline L$ consisting of the maps rational over $k$. The correspondence
$f\longrightarrow(f^{j}_{\alpha_1,\ldots,\alpha_r})$ is linear and one-to-one between $\underline L$ and the set all families of
elements in $K$ with only a finite number of nonzero components; since $\underline L_k$ is
mapped in this way on the set of all families with components in $k$, one sees that
$\underline L_k$ is a $k$-structure on $\underline L$. Furthermore, for any field $k'$, the set $\underline L_{k'}$ of the
multilinear maps rational over $k'$ is equal to $k'\unsure{\circ|\cdot}\underline L_k$.

Let finally $V$ and $W$ be two vector spaces with respective $k$-structures
$V_k$ and $W_k$, and let $f$ be a linear map from $V$ to $W$ rational over $k$. If a subspace $S$ of $V$ is rational over $k$, it is generated by $S\handwritten{\cap}V_k$ and therefore $f(S)$,
being generated by $f(S\handwritten{\cap}V_k)\handwritten{\subset}W_k$, is rational over $k$. Let now $T$ be a subspace
of $W$ rational over $k$; for a $k$-basis $\{c_\alpha\}$ of $K$, any vector $v$ in $V$ is uniquely expressed as $\sum_\alpha c_\alpha.v_\alpha$ for some $v_\alpha$ in $V_k$, and such a vector is in $T$ if and only if
$f(v_\alpha)$ is in $T$ for all $\alpha$ since $f(v)=\sum_\alpha c_\alpha.f(v_\alpha)$; therefore $f^{-1}(T)$ is generated by
$f^{-1}(T)\handwritten{\cap}V_k$ and is rational over $k$.

From the previous results, we shall deduce that the linear map $f$ is
rational over $k$ if and only if its graph $\Gamma$ is a subspace of $V\times W$ rational over $k$
(for the $k$-structure $V_k\times W_k$ on $V\times W$). Indeed, $\Gamma$ is the image of the map
$g:v\longrightarrow(v,f(v))$ of $V$ into $V\times W$ and if $f$ is rational over $k$, so is $g$ and so is
its image $\Gamma$. Conversely, assume $\Gamma$ rational over $k$; there exists therefore a
basis of $\Gamma$ consisting of vectors of $V\times W$ rational over $k$; such a basis is of
the form $\{(e_i,f(e_i))\}$ where $e_i$ and $f(e_i)$ are rational over $k$ for all $i$; but since
the projection from $\Gamma$ to $V$ is an isomorphism of vector spaces, the $e_i$ form a
basis of $V$ rational over $k$, and since the $f(e_i)$ are rational over $k$, so is $f$.

5. \uline{Conjugates}.

\uline{Let} $V$ \uline{be} \uline{a} \uline{vector} \uline{space} \uline{with} \uline{a} $k$-\uline{structure} $V_k$ \uline{and} \uline{let} $\sigma$ \uline{be} \uline{a} $k$-\uline{isomorphism}
\uline{from} \uline{a} \uline{field} $k'$ \uline{onto} \uline{a} \uline{field} $k''$.

We shall study the behavior under $\sigma$ of the vectors, the subspaces and the
multilinear maps. To begin with, I contend that there exists a unique map
$v\longrightarrow v^\sigma$ from $V_{k'}$ onto $V_{k''}$ such that:
% archive sheet number: 131
\scanpage{133}
\deleted{\unsure{(i)}}

$S$ \unsure{préschém.|préschéma} loc \unsure{noeth.}

$F$ foncteur \quad $(\mathrm{Sch})_{/S}\longrightarrow(\mathrm{Ens})$

(i) $F$ compatible avec limites inductives d'anneaux

(ii) $F$ compatible avec descente fidèlement plate \unsure{quasi cpcte.|quasi-compactes} \inserted{et localisation}

$\bigl[$(iii) Si $A$ anneau local noeth. \uline{complet} sur $S$,
corps résiduel $k$, alors $F(A)\overset{\sim}{\longrightarrow}F(k)$,

(iv) $F$ séparé.

(v) Pour \unsure{tt} corps $k$ sur $S$, $F(k)$ est \uline{fini},
et si $\unsure{K'|k'}$ est une ext. alg. close de $k$, \deleted{\unsure{$F(k$}} alg. clos,
$F(k)\unsure{-|\to}F(k')$ est bijectif.

Alors $F$ représentable par un \deleted{foncteur} objet
étale \deleted{\unsure{et qu.|et gpe}} séparé \unsure{qu.|quasi} fini sur $S$.

\deleted{\unsure{Réduction au cas|Réduisons-nous|Démonstration.}}

Supposons $S$ le spectre d'un anneau local complet
$A$, \deleted{et que} \inserted{à corps résiduel $k$ alg. clos.} de dim \unsure{$n$}, et qu'on sache
que sur $S_1=S-y$, le \deleted{foncteur} $F$ soit représentable
par $M_1$. \deleted{\unsure{Considérons les \illegible{1}}} Soient
$\unsure{\xi|\zeta}_i^0$ \deleted{\unsure{$\in F$|$\in k$}} \inserted{$(i\in I)$} les éléments de $\deleted{\unsure{R}}F(k)$, ils proviennent d'éléments
$\xi_i\in F(A)$ bien déterminés, considérons leurs
restrictions $\xi'_i=\xi_i\mid S_1\in F(S_1)$, ils correspondent à des
% archive sheet number: 132
\scanpage{135}
\unsure{1|I} \uline{Lemme} Soit $S$ un préschéma loc noeth, et
pour \deleted{tout} $s\in S$, soit $\deleted{\unsure{X_s}}$ $X(s)$ un préschéma
quasi fini \uline{étale} \uline{séparé} sur $\mathrm{Spec}(\mathcal O_s)$, et pour
$t\in S$, $s\in\overline t$, soit $\varphi_{st}:X(t)\unsure{-|\to}X(s)$ rendant
\diagram{\begin{tikzcd}[ampersand replacement=\&,row sep=small]
X(s)\arrow[d] \& X(t)\arrow[l]\arrow[d]\\
\mathrm{Spec}(\mathcal O_s) \& \mathrm{Spec}(\mathcal O_t)\arrow[l]
\end{tikzcd}}
cartésien le diagramme ci-contre. On suppose qu'on a transitivité
$\varphi_{st}\varphi_{tu}=\varphi_{su}$.
% margin note written sideways, linked to the end of the transitivity condition
\inserted{[On suppose que pour tt $s\in S$, il existe une \unsure{réalisation} étale \deleted{\illegible{1}}
$\mathfrak X$ de $X(s)$ sur un voisinage ouvert $U$ de $s$, \deleted{\unsure{tel que pour tt}} et pour tt $t\in U$,
un \unsure{hom} $\deleted{\unsure{\mathfrak X_s}}\mathfrak X\deleted{\unsure{s}}\times_U\mathrm{Spec}(\mathcal O_t)\to X(t)$, ces hom étant compatibles
avec les \unsure{\ldots}]}
% margin note written sideways, diagram alongside the inserted condition
\diagram{\begin{tikzcd}[ampersand replacement=\&,column sep=small,row sep=small]
X(s)\arrow[r]\arrow[d] \& \mathfrak X\arrow[d] \& \mathfrak X\times_U\mathrm{Spec}(\mathcal O_t)\arrow[l]\arrow[r]\arrow[d] \& X(t)\\
\mathrm{Spec}(\mathcal O_s)\arrow[r] \& U \& \mathrm{Spec}(\mathcal O_t)\arrow[l] \&
\end{tikzcd}}

\deleted{\unsure{des}} Alors il existe un préschéma
étale \inserted{quasi fini séparé} sur $S$, soit $X$, et des \inserted{$\mathcal O_s$-}isom
\[
\varphi_s:X\times_S\mathrm{Spec}(\mathcal O_s)\unsure{-|\to|\overset{\sim}{\longrightarrow}}X(s)
\]
rendant commutatifs les diagrammes qu'on pense.
$X$ est unique à isom unique près.

Ce dernier \deleted{fait} résulte d'un
énoncé plus général sur la catégorie des
préschémas de type fini sur $S$, qui donne
un foncteur \unsure{pleinement} fidèle. Il nous permet
de nous ramener au local et supposer $S$
affine, donc noeth. On procède par
récurrence noethérienne, en supposant l'énoncé
prouvé pour les sous-préschémas \unsure{fermés} de
$S\ne S$.
% archive sheet number: 134
\scanpage{137}
Soit $S$ un schéma \uline{réduit} \uline{connexe}.

On considère des schémas en groupes, finis et plats sur $S$,
soit $G$, pour lesquels on a une suite exacte
\[
e\unsure{-|\to}\deleted{\unsure{H}}U\to G\to\Gamma\to e
\]
où \deleted{\unsure{le}}la projection $U\to S$ est plate, finie, \deleted{\unsure{\illegible{1}|et}} et
géométriquement \deleted{\unsure{de rang \illegible{1} supérieur|homomorphisme}} injectif,
et où $\Gamma\unsure{-|\to}S$ est plat, fini, séparable. Alors le
sous-$S$-groupe $U$ de $G$ est déterminé de
façon unique [pour cette assertion, l'hypothèse
``$S$ réduit'' est-elle nécessaire ?].
% slanted
\margin{$S$-\uline{Groupes} \uline{admissibles} ($S$ réduit ?)}
On veut étudier des fibrés principaux \inserted{\unsure{hom.}} $P$ sous
de tels $S$-groupes. On suppose $S$ ponctué
par $a$\deleted{\unsure{$\xi$}}, et on considère des $P$ avec
structure ponctuée.

\uline{Proposition} \deleted{\unsure{soit \illegible{1}}} Soit $S'\to S$ un morphisme
\deleted{\unsure{plat fini surjectif fidèlement plat}}, $G$ un $S$-groupe
fini.
% slanted
\margin{$S'$ et $S$ réduits}
\deleted{\unsure{Soit $G\times_S S'$ est admissible, ?}} \inserted{Si $G\deleted{\unsure{\times_S S'}}$ est admissible,} $G\times_S S'$ aussi,
la réciproque est vraie si $S'\to S$ est fidèlement plat.

Direct est trivial. Réciproque résulte \unsure{de|du|des}
le théorème de \unsure{la descente} et unicité.
% archive sheet number: 136
\scanpage{138}
\uline{Corollaire} Supposons $\deleted{\unsure{P}}G$ admissible, et $P$ un
fibré principal sous $G$. Alors $\mathrm{ad}(P)$, considéré
comme $S$-groupe, est admissible.

\uline{Proposition 2} Soit $(P,\xi)$ un fibré principal ponctué
sous un $S$-groupe admissible $G$. Alors
$\mathrm{Aut}(P,\xi)$ est réduit à l'identité.

Cela résulte, moyennant le corollaire précédent,
\deleted{de}

\deleted{\uuline{Lemme}} \uline{Proposition 3} Deux sections d'un
$P$ (\inserted{principal sous $G$} admissible sur $S$ \inserted{réduit connexe}) qui ont $=$ valeur
en un pt, sont identiques. [De façon
plus précise, si $s$ est une section de $P/S$,
alors il existe une composante connexe
$\deleted{\unsure{P_0}}Q$ de $P$, tel que $Q_{\text{réd}}\to S$ soit un
isom, et l'isomorphisme réciproque induit par $s$]

Un $(P,\xi)$ \deleted{est dit} (relatif à $G$ admissible)
est dit \uline{minimal} si on \unsure{ne peut trouver|ne trouve pas}
de \unsure{$S$-ss-groupe|$S$-sous-groupe|sous-groupe} admissible $G'$ de $G$, et
$(P',\xi')$ relatif à $G$, t.q. $(P,\xi\deleted{\unsure{'}})\simeq(P',\xi')\times_G G'$.
% archive sheet number: 137
\scanpage{139}
\deleted{\uline{Proposition 4}} Soit $\unsure{\mathcal Z^{1}|\mathcal Z^{0}|Z^{1}|\mathfrak Z^{1}|\mathfrak Z^{*}}(a;G)=$ ens. des classes à isom
près de $G$-fibrés principaux hom. \uline{\unsure{ponctués|pointés}} au
dessus de $a$. C'est un foncteur covariant en
le $S$-groupe admissible $G$.

\uline{Proposition 4} Si $\alpha\in\mathcal Z^{1}(a;G)$ est \uline{minimal},
alors $u\longmapsto u_{*}(\alpha)$ de $\operatorname{Hom}^{\mathrm{gr}.}_{S}(G,G')$ dans
$\deleted{\unsure{Z|\mathfrak Z}}\mathcal Z(a;G')$ est \uline{injective}.
\margin{\uline{\unsure{Sans rigueur|Sens rigoureux}}} % slanted

Soient en effet $u,v:G\mathbin{-}G'$ t.q.
$u\alpha=v\alpha$. Soit $(P,\xi)$ de classe $\alpha$,
$(P',\xi')$ de classe $u\alpha=v\alpha$, $(u,v):G\to G'\times G'$ \deleted{\unsure{sur}}
(produits \uline{sur $S$}) d'où $\overbrace{(P,\xi)\mathbin{\underset{G}{\times}}(G'\times G')}^{(Q,\eta)}$ ;
par hypothèse, il y a un \unsure{hom|mor.}
(compatible avec les ponctuations, et avec \inserted{la diagonale}
$G'\mathbin{-}G'\times G'$) $(P',\xi')\longrightarrow(P,\xi)\mathbin{\underset{G}{\times}}(G'\times G')$, d'où
un diagramme
\diagram{\begin{tikzcd}[ampersand replacement=\&]
(P,\xi)\arrow[r] \& (\deleted{\unsure{P}}Q,\eta)\\
\& (P',\xi')\arrow[u]
\end{tikzcd}}
au dessus de
\diagram{\begin{tikzcd}[ampersand replacement=\&]
G\arrow[r,"{(u,v)}"] \& G'\times G'\\
\& G'\arrow[u]
\end{tikzcd}}
% archive sheet number: 138
\scanpage{140}
Soit $(R,\lambda)$ le produit fibré de
$(P,\xi)$ et $(P',\xi')$ sur $(Q,\eta)$,
% archive sheet number: 139
